% SGA 3，第 VII 篇_B，英文重建本——第 1 轮（中文译本）。
% 范围：依据原文件第 20--23 页，包含 1.2.3 全部；止于原文件第 23 页的 1.2.4 之前。
% 合并阅读版页码： 533--536. 再版印刷页码： 522--525.
% 本范围内的原版页边定位： 495--497.
% 依据： Polo--Gille Exp7B-13oct24.pdf,
% SHA-256 D303690E0195654BBE7C3DAA5D89EA23E24CD709A3575217234AC19AE8836694.
% 曾参阅 Jacob Reinhold 的 CC BY 4.0 Markdown，作为对比版本来源。
% 本范围内无原生图；所有陈列数学均为原生 LaTeX。

\medskip
\noindent\textbf{1.2.3.}\footnote{\label{sga3:VIIB:nde:45}编者添加了下列引理，并引入编号 \SGARefLink{sga3:VIIB:1.2.3.A}{1.2.3.A} 至
\SGARefLink{sga3:VIIB:1.2.3.F}{1.2.3.F}，以供下文引用。}
\phantomsection
\label{sga3:VIIB:1.2.3}
先从下列引理开始。

\medskip
\noindent\textbf{引理 1.2.3.A。}
\phantomsection
\label{sga3:VIIB:1.2.3.A}
{\itshape 设 \(k\longrightarrow K\) 是伪紧环态射，\(B\) 是伪紧 \(K\)-模，
\(M\) 是不带拓扑的投射 \(k\)-模。存在伪紧 \(K\)-模的典范同构
\begin{equation}
  \theta:
  \operatorname{Hom}_k(M,k)\widehat{\otimes}_k B
  \xrightarrow{\ \sim\ }
  \operatorname{Hom}_k(M,B).
  \tag{2}
  \label{sga3:VIIB:1.2.3:eq2}
\end{equation}}

这里 \(M^*=\operatorname{Hom}_k(M,k)\) 赋予
\SGARefLink{sga3:VIIB:0.2.2}{0.2.2} 中定义的拓扑，使之成为伪紧 \(k\)-模。
\(\operatorname{Hom}_k(M,B)\) 上的拓扑类似定义：\(0\) 的邻域基由以下
\(K\)-子模构成，其中 \(x\in M\)，\(B'\) 是 \(B\) 的开 \(K\)-子模：
\[
  \mathscr H(x,B')
  =
  \{f\in\operatorname{Hom}_k(M,B)\mid f(x)\in B'\}.
\]
最后，\(M^*\widehat{\otimes}_kB\) 是从 \(M^*\) 通过基变换得到的伪紧
\(K\)-模；参见\ \SGARefLink{sga3:VIIB:0.5.A}{0.5.A}。

采用这些约定，\(M=k\) 时 \(\theta\) 显然是同构。此外，将 \textup{(2)} 的两端
视为关于 \(M\) 的函子，它们都将直和变为积；特别地，它们与有限直和交换。
因此，先得到 \(M\) 为自由 \(k\)-模时 \textup{\SGARefLink{sga3:VIIB:1.2.3:eq2}{(2)}}
是同构，再得到 \(M\) 为投射模时也是如此。

\medskip
\noindent\textbf{定义 1.2.3.B。}
\phantomsection
\label{sga3:VIIB:1.2.3.B}
设 \(N\) 是伪紧 \(k\)-模。我们用 \(\mathbf V_k^f(N)\) 或
\(\mathbf N^\dagger\)\footnote{\label{sga3:VIIB:nde:46}事实上，第二种记号不会使用；
见编者注~\SGARefLink{sga3:VIIB:nde:47}{(47)}。}表示将每个
\(C\in\operatorname{Ob}(\operatorname{Alf}/k)\) 对应到如下 \(C\)-模的 \(\mathbf O_k\)-模
\[
  (C\widehat{\otimes}_kN)^\dagger,
\]
这是 \(C\widehat{\otimes}_kN\) 在 \SGARefLink{sga3:VIIB:0.2.2}{0.2.2} 意义下的对偶；
即 \(C\)-模
\[
  \operatorname{Hom}_c(C\widehat{\otimes}_kN,C)
  \simeq
  \operatorname{Hom}_c(N,C|_k),
\]
其中 \(C|_k\) 表示从 \(C\) 通过限制标量得到的 \(k\)-模。
这个 \(\mathbf O_k\)-模 \(\mathbf V_k^f(N)\) 称为 \(N\) 的\emph{对偶}。

若 \(k_{\mathfrak m}\) 是 \(k\) 的局部分量，则对
\(\operatorname{Alf}/k\) 的每个对象 \(C\)，
\[
  \operatorname{Hom}_c(k_{\mathfrak m},C|_k)
  =
  C_{\mathfrak m}
  =
  C\otimes_k k_{\mathfrak m}
\]
是投射 \(C\)-模。此外，
\[
  D\otimes_C C_{\mathfrak m}=D_{\mathfrak m}
\]
对 \(\operatorname{Alf}/k\) 中的每个态射 \(C\longrightarrow D\) 成立。
因此，\(\mathbf V_k^f(k_{\mathfrak m})\) 是平坦 \(\mathbf O_k\)-模；
参见\ \SGARefLink{sga3:VIIB:1.2.1}{1.2.1}。由于每个投射伪紧 \(k\)-模都是
\(k_{\mathfrak m}\) 的积（\SGARefLink{sga3:VIIB:0.2.1}{命题~0.2.1}），
\SGARefLink{sga3:VIIB:0.2.F}{推论~0.2.F} 给出：

\medskip
\noindent\textbf{推论 1.2.3.C。}
\phantomsection
\label{sga3:VIIB:1.2.3.C}
\textit{只要 \(N\) 是投射伪紧 \(k\)-模，\(\mathbf O_k\)-模
\(\mathbf V_k^f(N)\) 就是平坦的。}

\medskip
\noindent\textbf{定义 1.2.3.D。}
\phantomsection
\label{sga3:VIIB:1.2.3.D}
反之，设 \(\mathbf M\) 是容许 \(\mathbf O_k\)-模
（参见\ \SGARefLink{sga3:VIIB:1.2.1}{1.2.1}）。\footnote{
\label{sga3:VIIB:nde:47}除非假定 \(\mathbf M\) 容许，否则编者无法理解下面的定义，
因此加入了这一假设。编者还以 \(\Gamma^*(\mathbf M)\) 代替 \(\mathbf M^*\)，
以免将函子 \(\mathbf M\) 的对偶模 \(\mathbf M^*\) 与模 \(N\) 的对偶函子
\(\mathbf N^\dagger\) 混淆；参见\ 编者注~(46)。最后，编者展开了
\(\Gamma^*(\mathbf M)\) 的构造。}
我们将下列伪紧 \(k\)-模称为 \(\mathbf M\) 的\emph{对偶}，并记作
\(\Gamma^*(\mathbf M)\)。令 \(\mathfrak l\) 遍历 \(k\) 的开理想，赋予每个
\(k/\mathfrak l\)-模
\[
  \mathbf M(k/\mathfrak l)^*
  =
  \operatorname{Hom}_{k/\mathfrak l}
    \bigl(\mathbf M(k/\mathfrak l),k/\mathfrak l\bigr)
\]
以 \SGARefLink{sga3:VIIB:0.2.2}{0.2.2} 中描述的拓扑；于是它是伪紧 \(k\)-模。
由于
\[
  \mathbf M(k/\mathfrak l)
  \simeq
  \mathbf M(k/\mathfrak l')
    \otimes_{k/\mathfrak l'}(k/\mathfrak l)
  \qquad
  (\mathfrak l\supset\mathfrak l'),
\]
存在过渡态射
\[
  \begin{split}
  \operatorname{Hom}_{k/\mathfrak l'}
    \bigl(\mathbf M(k/\mathfrak l'),k/\mathfrak l'\bigr)
  &\longrightarrow
  \operatorname{Hom}_{k/\mathfrak l'}
    \bigl(\mathbf M(k/\mathfrak l'),k/\mathfrak l\bigr)\\
  &=
  \operatorname{Hom}_{k/\mathfrak l}
    \bigl(\mathbf M(k/\mathfrak l),k/\mathfrak l\bigr).
  \end{split}
\]
按定义，\(\Gamma^*(\mathbf M)\) 就是这个伪紧 \(k\)-模的滤过投射系的投射极限。

从现在起，还假定 \(\mathbf M\) 是平坦 \(\mathbf O_k\)-模
（参见\ \SGARefLink{sga3:VIIB:1.2.1}{1.2.1}）。于是每个
\(\mathbf M(k/\mathfrak l')\) 都是投射 \(k/\mathfrak l'\)-模，
所以上述过渡态射是满射。因此，投影
\[
  \Gamma^*(\mathbf M)
  \longrightarrow
  \mathbf M(k/\mathfrak l)^*
\]
也是满射，因为 \(\operatorname{PC}(k)\) 中的滤过投射极限正合。

设 \(\tau:k\longrightarrow K\) 是伪紧环态射。
我们用 \(\mathbf M\otimes_kK\)，或简记 \(\mathbf M_K\)，表示如下定义的
\(\mathbf O_K\)-函子。若 \(C\) 是有限长度的 \(K\)-代数，
则 \(k\longrightarrow C\) 的核是开理想 \(\mathfrak l\)，我们令
\[
  \mathbf M_K(C)
  =
  \mathbf M(k/\mathfrak l)\otimes_k C.
\]
对于每个开理想 \(\mathfrak l'\subset\mathfrak l\)，
\(\mathbf M(k/\mathfrak l')\otimes_kC\) 给出同样的值。
然后将 \(\Gamma_K^*(\mathbf M_K)\) 定义为下列伪紧 \(K\)-模的投射极限，
其中 \(I\) 遍历 \(K\) 的开理想：
\[
  \operatorname{Hom}_{K/I}\bigl(\mathbf M_K(K/I),K/I\bigr)
  =
  \operatorname{Hom}_{k/\mathfrak l}
    \bigl(\mathbf M(k/\mathfrak l),K/I\bigr),
\]
其中右端的 \(\mathfrak l\) 是 \(k\) 的任意满足 \(\tau(\mathfrak l)\subset I\) 的开理想。

由\SGARefLink{sga3:VIIB:1.2.3.A}{引理~1.2.3.A}，右端等同于
\[
  \mathbf M(k/\mathfrak l)^*
  \widehat{\otimes}_k(K/I).
\]
由于投影 \(\Gamma^*(\mathbf M)\longrightarrow\mathbf M(k/\mathfrak l)^*\)
是满射，后面这些模的投射极限恰是伪紧 \(K\)-模
\[
  \Gamma^*(\mathbf M)\widehat{\otimes}_kK
\]
（参见\ \SGARefLink{sga3:VIIB:0.5.A}{0.5.A}）。因此已经证明：对每个平坦
\(\mathbf O_k\)-模 \(\mathbf M\)，\(\Gamma^*(\mathbf M)\) 的形成与基扩张交换；即
\begin{equation}
  \Gamma_K^*(\mathbf M\otimes_kK)
  \simeq
  \Gamma^*(\mathbf M)\widehat{\otimes}_kK.
  \tag{\(\star\)}
  \label{sga3:VIIB:1.2.3:star}
\end{equation}

\medskip
\noindent\textbf{命题 1.2.3.E。}
\phantomsection
\label{sga3:VIIB:1.2.3.E}
{\itshape
\begin{enumerate}
  \item[\textup{(i)}]
  函子
  \[
    N\longmapsto\mathbf V_k^f(N),
    \qquad
    \mathbf M\longmapsto\Gamma^*(\mathbf M)
  \]
  在投射伪紧 \(k\)-模范畴与平坦 \(\mathbf O_k\)-模范畴之间定义反等价。
  \footnote{\label{sga3:VIIB:nde:48}编者添加了下面的断言
  \textup{\SGARefLink{sga3:vii:paragraph:5-1:equation:ii}{(ii)}} 及随后的证明。
  原文说：“若 \(\mathbf M\) 是平坦 \(\mathbf O_k\)-模，显然 \(\mathbf M\)
  是 \(\mathbf M^*\) 的对偶；因此函子 \(N\mapsto N^*\) 和
  \(\mathbf M\mapsto\mathbf M^*\) 定义反等价 \ldots”}

  \item[\textup{(ii)}]
  此外，若 \(k\longrightarrow K\) 是伪紧环态射，则上述反等价在以下意义下
  “与基变换交换”：若 \(N\simeq\Gamma^*(\mathbf M)\)，则
  \[
    N\widehat{\otimes}_kK
    \simeq
    \Gamma_K^*(\mathbf M\otimes_kK).
  \]
\end{enumerate}}

\noindent\textit{证明。}
首先，有自然变换
\[
  \phi_N:
  N\longrightarrow\Gamma^*\bigl(\mathbf V_k^f(N)\bigr).
\]
由 \SGARefLink{sga3:VIIB:0.3.5}{0.3.5} 和 \SGARefLink{sga3:VIIB:0.2.F}{0.2.F}，
函子 \(\Gamma^*\circ\mathbf V_k^f\) 与积交换，所以只需验证 \(N\) 为
\(k\) 的局部分量 \(k_{\mathfrak m}\) 时，\(\phi_N\) 是同构。
在这种情形下，对 \(k\) 中包含于 \(\mathfrak m\) 的每个开理想 \(\mathfrak l\)，自然态射
\[
  (k/\mathfrak l)_{\mathfrak m}
  \longrightarrow
  \operatorname{Hom}_{k/\mathfrak l}
  \left(
    \operatorname{Hom}_c
      (k_{\mathfrak m}\widehat{\otimes}_k k/\mathfrak l,
       k/\mathfrak l),
    k/\mathfrak l
  \right)
\]
是同构，这就证明了该断言。

反之，设 \(\mathbf M\) 是平坦 \(\mathbf O_k\)-模。下面证明
\[
  \Gamma^*(\mathbf M)
  =
  \varprojlim_{\mathfrak l}\mathbf M(k/\mathfrak l)^*
\]
是 \(\operatorname{PC}(k)\) 的投射对象。设 \(N\longrightarrow N'\) 是
\(\operatorname{LF}(k)\) 中对象之间的满态射。由
\SGARefLink{sga3:VIIB:0.2.F}{0.2.F}\,\textup{\SGARefLink{sga3:vii:paragraph:5-1:equation:i}{(i)}} 和
\textup{\SGARefLink{sga3:vii:paragraph:5-1:equation:ii}{(ii)}}，只需证明自然映射
\[
  \varinjlim_{\mathfrak l}
    \operatorname{Hom}_c(\mathbf M(k/\mathfrak l)^*,N)
  \longrightarrow
  \varinjlim_{\mathfrak l}
    \operatorname{Hom}_c(\mathbf M(k/\mathfrak l)^*,N')
\]
是满射。这是显然的：对某个开理想 \(\mathfrak l_0\)，\(N\) 和 \(N'\) 都是
\(k/\mathfrak l_0\)-模；若 \(\mathfrak l\subset\mathfrak l_0\)，
每个态射 \(\mathbf M(k/\mathfrak l)^*\longrightarrow N'\) 都提升为态射
\(\mathbf M(k/\mathfrak l)^*\longrightarrow N\)，因为
\(\mathbf M(k/\mathfrak l)^*\) 是 \(\operatorname{PC}(k/\mathfrak l)\) 的投射对象。

存在从 \(\operatorname{Alf}/k\) 到 \((\mathbf{Sets})\) 的函子之间的态射
\[
  \psi:
  \mathbf M\longrightarrow
  \mathbf V_k^f\bigl(\Gamma^*(\mathbf M)\bigr)
\]
设 \(B\) 是 \(\operatorname{Alf}/k\) 的对象。我们断言
\begin{equation}
  \begin{split}
  \psi(B):\mathbf M(B)
  &\longrightarrow
  \mathbf V_k^f\bigl(\Gamma^*(\mathbf M)\bigr)(B)\\
  &=
  \operatorname{Hom}_c
  \left(
    \varprojlim_{\mathfrak l}
      \bigl(\mathbf M(k/\mathfrak l)^*
        \widehat{\otimes}_kB\bigr),
    B
  \right)
  \end{split}
  \tag{1}
  \label{sga3:VIIB:1.2.3:eq1}
\end{equation}
是双射。在上面的等式中，我们使用了 \(\widehat{\otimes}\) 与滤过投射极限交换这一事实。
设 \(\mathfrak l_0\) 是 \(k\) 中包含于 \(k\longrightarrow B\) 的核的开理想。
由\SGARefLink{sga3:VIIB:1.2.3.A}{引理~1.2.3.A}，对每个
\(\mathfrak l\subset\mathfrak l_0\)，都有伪紧 \(B\)-模的典范同构
\[
  \mathbf M(k/\mathfrak l)^*\widehat{\otimes}_kB
  \xrightarrow{\ \sim\ }
  \operatorname{Hom}_{k/\mathfrak l}
    \bigl(\mathbf M(k/\mathfrak l),B\bigr).
\]
由于
\[
  \mathbf M(B)
  =
  \mathbf M(k/\mathfrak l)
    \otimes_{k/\mathfrak l}B,
\]
右端是 \(\operatorname{Hom}_B(\mathbf M(B),B)\)。因此，
\textup{\SGARefLink{sga3:VIIB:1.2.3:eq1}{(1)}} 中的投射系在
\(\mathfrak l\subset\mathfrak l_0\) 时为常值，而 \textup{(1)} 归结为典范态射
\[
  \mathbf M(B)
  \longrightarrow
  \operatorname{Hom}_c
  \bigl(\operatorname{Hom}_B(\mathbf M(B),B),B\bigr).
\]
由 \SGARefLink{sga3:VIIB:0.2.2}{0.2.2}，这是同构，因为 \(B\) 是 Artin 环，
而 \(\mathbf M(B)\) 是投射 \(B\)-模。这证明了命题的断言
\textup{\SGARefLink{sga3:vii:paragraph:5-1:equation:i}{(i)}}；断言
\textup{\SGARefLink{sga3:vii:paragraph:5-1:equation:ii}{(ii)}} 由上面建立的同构
\textup{(\(\star\))} 得到。
\hfill\(\square\)

\medskip
\noindent\textbf{注记 1.2.3.F。}
\phantomsection
\label{sga3:VIIB:1.2.3.F}
回到命题的陈述，并进一步假定 \(N\) 是拓扑自由的伪紧 \(k\)-模。
这时，可以“协调一致地”为各 \(C\)-模 \(\mathbf V_k^f(N)(C)\) 选择基。

事实上，设 \((n_i)\) 是 \(N\) 的伪基（\SGARefLink{sga3:VIIB:0.2.1}{0.2.1}），
对 \(C\in\operatorname{Alf}/k\)，设 \(n_i^C\) 是 \(n_i\) 在
\(C\widehat{\otimes}_kN\) 中的典范像。按下式定义
\(\delta_i^C\in(C\widehat{\otimes}_kN)^\dagger\)：
\[
  \delta_i^C(n_i^C)=1,
  \qquad
  \delta_i^C(n_j^C)=0\quad(i\ne j).
\]
则 \((\delta_i^C)\) 是 \(\mathbf V_k^f(N)(C)\) 的基。此外，对
\(\operatorname{Alf}/k\) 中的每个态射 \(\varphi:C\longrightarrow D\)，
映射 \(\mathbf V_k^f(N)(\varphi)\) 将 \(\delta_i^C\) 映为 \(\delta_i^D\)。

% 下一断点： Exp7B 原文件页码 23, 合并阅读版页码 536,
% 再版印刷页码 525, 紧接其前： 1.2.4.
